\documentclass[11pt]{letter}
\usepackage[margin=2.5cm]{geometry}
\usepackage{mathpazo}
\usepackage{url}
\usepackage{hyperref}
\hypersetup{colorlinks=true,urlcolor=blue,linkcolor=blue}

\date{April 27, 2026}

\signature{Yan Ma \\ Lizhuo Zhang \\ Hunan Agricultural University \\ Changsha 410128, China \\ E-mail: zhanglizhuo@hunau.edu.cn}

\address{College of Education \\ Hunan Agricultural University \\ Changsha 410128, China}

\begin{document}

\begin{letter}{Editorial Office \\ \textit{PLOS ONE} \\ Public Library of Science}

\opening{Dear Editor,}

We would like to submit the enclosed manuscript, ``A Four-Dimension Pre-Modeling Audit Protocol for Educational Prediction Benchmarks: Evidence from Seven Public Datasets,'' for consideration as an original Research Article in \textit{PLOS ONE}.

\textbf{Background and motivation.}
Predictive models in education increasingly inform high-stakes decisions about students and teachers, yet the benchmarks that underpin these models are rarely stress-tested before modeling begins. In computer vision and natural language processing, benchmark-audit studies have become routine, but no comparable audit protocol exists for the educational prediction domain.

\textbf{Summary of contributions.}
This study makes three contributions. First, we propose a reusable four-dimension pre-modeling audit protocol covering baseline gap, split instability, null separation, and metadata adequacy (operationalized as group-aware holdout). Second, we validate the protocol on seven publicly available educational datasets spanning $N = 145$ to $32{,}593$, revealing that five of seven fail metadata adequacy: four fail linear-model group-aware holdout (two with model-invariant confounding, two with remediable linear-specific overfitting), and one cannot be tested because grouping metadata is unavailable. Third, a model-complexity ablation demonstrates that data-level fragility is algorithm-invariant: escalating from linear baselines to ensemble methods amplifies instability on fragile datasets rather than resolving it. A classification-metric sensitivity check confirms that audit conclusions are robust to metric choice.

\textbf{Relevance to \textit{PLOS ONE}.}
\textit{PLOS ONE} judges manuscripts on the technical soundness and rigor of the research rather than on perceived novelty or scope, and explicitly welcomes methodological and replication studies of broad relevance. The present work fits this remit precisely: it offers a transparent, reproducible quality-control protocol for an entire class of widely used educational datasets, applies it uniformly to seven public benchmarks, and reports both supporting and disconfirming evidence dataset by dataset. The full code, data sources, and figure-regeneration scripts are released so that the analyses can be re-executed by any reader.

\textbf{Compliance with \textit{PLOS ONE} criteria.}
The study is original and has not been published previously, nor is it under consideration elsewhere. All listed authors meet authorship criteria, have read and approved the manuscript, and have agreed to its submission. The work uses only publicly available, de-identified educational datasets and does not involve any new human-subjects data collection; institutional ethical review was therefore not required. The authors declare no competing interests. Funding was provided by the Education Department of Hunan Province, China (Project Number: HNJG-2022-0608); the funder had no role in study design, data analysis, decision to publish, or manuscript preparation. All seven audited datasets are publicly available from their original maintainers and are cited in the manuscript; the audit protocol implementation and all analysis and figure-generation scripts are available at \url{https://github.com/zhanglizhuo/BehaviorAudit}, satisfying the \textit{PLOS} Data Availability policy.

We believe the findings will be of broad interest to readers in educational data mining, learning analytics, educational measurement, and applied machine learning. We look forward to your consideration.

\closing{Sincerely,}

\end{letter}
\end{document}
